\section{The self-interaction: nine channels, four couplings}
\label{sec:selfint}

This section analyzes the quartic term $\Lself$ of Eq.~(\ref{eq:Lself}) by exact algebra. Its results
are used in every later section: the number of independent couplings, the couplings that act in
uniform matter, the couplings that act between two particles, and the symmetries that the
self-interaction keeps or breaks. \SIG{On the $4+4$ spacetime the quartic term can be completed in two
ways once the theory is quantized. Section~\ref{sec:forms} defines them: \emph{form C}, the covariant
form, which is the definition of the model in this paper, and \emph{form Q}, the quantized-sector
form, which is the alternative and is recorded beside it. The later subsections give the algebra for
both forms, and Sec.~\ref{sec:formsdisc} discusses the two.}{}

%%<1644>
\subsection{The quantized sector and the two forms of the self-interaction}
\label{sec:forms}

On the $4+4$ spacetime only fields that do not depend on the three additional times $x5,x6,x7$ have a
unitary evolution (Sec.~\ref{sec:dirac}); they form the \emph{quantized sector}. For such fields the
kinetic and mass terms (\ref{eq:Lfree}) become, with $\beta=-i\gamma^{4}$,
$\alpha^{j}=-\gamma^{4}\gamma^{j}$ ($j=0,1,2,3$), the Hamiltonian matrix
$\hh=\sum_j\alpha^{j}(-i\partial_j)+\beta m$, the matrix $\Wg=i\,\Cm\gamma^{4}$ and the identity
$\Cm=-\Wg\beta$,
\begin{equation}
\tilde\psi\big(\gamma^{A}\partial_A-m\big)\psi=-\psi^\dagger\,\Wg\,\big(i\partial_4-\hh\big)\psi .
\label{eq:indef}
\end{equation}
The matrix $\Wg$ is Hermitian with $\Wg^{2}=1$; it has the eigenvalues $+1$ and $-1$, eight times
each, and it commutes with $\hh$. The field therefore consists of two \emph{halves},
\begin{equation}
\psi=\psi_++\psi_-,\qquad\Wg\,\psi_\pm=\pm\,\psi_\pm ,
\label{eq:halves}
\end{equation}
and $\hh$ acts inside each half: of the eight states of positive energy of a given momentum, four lie
in each half. By Eq.~(\ref{eq:indef}) the covariant Lagrangian is the ordinary Lagrangian
$\psi^\dagger(i\partial_4-\hh)\psi$ of a Dirac field on the half $\Wg=-1$, and \emph{minus} that
Lagrangian on the half $\Wg=+1$. For a free field the sign does not change the equation of motion, and
the quantum theory is defined by the positive Lagrangian
\begin{equation}
\begin{aligned}
\mathcal L_{\mathrm{kin}}+\mathcal L_{\mathrm{mass}}&=\psi^\dagger\big(i\partial_4-\hh\big)\psi\\
&=\bar\psi\,(i\beta\partial_4-\beta\hh)\,\psi,\qquad\bar\psi=\psi^\dagger\beta,
\end{aligned}
\label{eq:Lpos}
\end{equation}
for all sixteen components, which gives a particle eight internal states of positive norm. Its mass
term is $-m\,\bar\psi\psi$, the same mass for both halves. The mass term of the covariant Lagrangian
is different:
\begin{equation}
-m\,\tilde\psi\psi=m\,\big(\psi_+^\dagger\beta\,\psi_+-\psi_-^\dagger\beta\,\psi_-\big),
\label{eq:covmass}
\end{equation}
a mass term of opposite sign in the two halves. For a free field this difference is invisible,
because the spectrum of $\hh$ does not depend on the sign of $m$. With a self-interaction the two
halves are no longer equivalent, and the quantized theory can be completed in two ways. Both keep
the positive Lagrangian (\ref{eq:Lpos}); they differ in the adjoint with which the nine invariants
are built.

\emph{Form C, the covariant form: the definition.} The quartic term is that of the covariant
Lagrangian, Eq.~(\ref{eq:Lself}), with the adjoint $\tilde\psi=\psi^\dagger\Cm$ and the matrices
(\ref{eq:Theta}):
\begin{equation}
\begin{aligned}
\mathcal L_{\mathrm{kin}}+\mathcal L_{\mathrm{mass}}+\LselfC&=\psi^\dagger\big(i\partial_4-\hh\big)\psi-\frac12\sum_{r=0}^{8}g_r\,I_r^{\mathrm C},\\
I_r^{\mathrm C}&=\sum_{|K|=r}\sgnK\,\big(\tilde\psi\,\Theta^{K}\psi\big)^{2} .
\end{aligned}
\label{eq:LselfC}
\end{equation}
($I_r^{\mathrm C}$ is the invariant $I_r$ of Eq.~(\ref{eq:Lself}); the superscript distinguishes it
from the invariant of the second form.) \emph{This is the definition of the $4+4$ model in this
paper.} Its quartic term is invariant under all 28 rotations and boosts of the frame, as the
hypothesis requires of the invariants of the spacetime. Its kinetic and mass terms are invariant
under the 13 of them that commute with $\Wg$: the ten rotations and boosts of $x0,\dots,x4$ and the
three rotations of $x5,x6,x7$. The other fifteen mix one of $x0,\dots,x4$ with an additional time
and do not preserve the quantized sector. The field equation is
\begin{equation}
i\partial_4\psi=\hh\,\psi+\sum_rg_r\sum_{|K|=r}\big(\tilde\psi\Theta_K\psi\big)\,\Cm\Theta^{K}\psi ;
\label{eq:fieldC}
\end{equation}
it conserves $\int\psi^\dagger\psi$, the number of particles of both halves together.

\emph{Form Q, the quantized-sector form: the alternative.} The bilinears are formed with the adjoint
$\bar\psi=\psi^\dagger\beta$ of Eq.~(\ref{eq:Lpos}) and with the matrices
\begin{equation}
\gamma_{44q}^{A}=-i\gamma^{A}\ (A=0,\dots,4),\qquad\gamma_{44q}^{A}=\gamma^{A}\ (A=5,6,7),
\label{eq:gq}
\end{equation}
which satisfy $\gamma_{44q}^{4}=\beta$ and in terms of which the positive Lagrangian reads
$\bar\psi(i\gamma_{44q}^{A}\partial_A-m)\psi$ for fields of the quantized sector. They have the
squares of the Dirac matrices of a spacetime with one time: $+1$ for $A=4$ and $-1$ otherwise. With
\begin{equation}
\begin{aligned}
\ThQ^{K}&=\zeta_r^{Q}\,\gamma_{44q}^{A_1}\cdots\gamma_{44q}^{A_r},\\
\zeta_r^{Q}&=\begin{cases}i,&r=2,3,6,7,\\1,&\text{otherwise},\end{cases}
\end{aligned}
\label{eq:ThetaQ}
\end{equation}
and with the sign $\sgnQ=\pm1$, the product of the squares of the matrices $\gamma_{44q}^{A}$ in $K$,
so that $(\ThQ^{K})^{2}=\sgnQ$, the invariants and the self-interaction of this form are
\begin{equation}
I_r^{\mathrm Q}=\sum_{|K|=r}\sgnQ\,\big(\bar\psi\,\ThQ^{K}\psi\big)^{2},\qquad\LselfQ=-\frac12\sum_rg_r\,I_r^{\mathrm Q} .
\label{eq:LselfQ}
\end{equation}
The field equation,
\begin{equation}
i\partial_4\psi=\hh\,\psi+\sum_rg_r\sum_{|K|=r}\sgnQ\,\big(\bar\psi\,\ThQ^{K}\psi\big)\,\beta\,\ThQ^{K}\psi ,
\label{eq:fieldQ}
\end{equation}
is term by term that of the $7+1$ spacetime for fields that do not depend on $x5,x6,x7$, because the
algebra of the matrices (\ref{eq:gq}) is that of the $7+1$ Dirac matrices: every result for such
fields on $7+1$ holds for form Q without change. Its invariants are invariant under the group that
the 28 products $\tfrac12\gamma_{44q}^{A}\gamma_{44q}^{B}$ generate. Thirteen of these are the
transformations of the coordinates named above; the other fifteen mix the two halves of the field
and are not transformations of the coordinates of the $4+4$ spacetime.

\emph{The exact relation between the two forms.} For every one of the 256 lists $K$,
$\Cm\Theta^{K}=z_K\,\Wg\,\beta\,\ThQ^{K}$ with a phase $z_K=\pm1$ if $K$ contains an even number of the
indices 5, 6, 7 and $z_K=\pm i$ if it contains an odd number. Consequently
\begin{equation}
I_r^{\mathrm C}=\sum_{|K|=r}\sgnQ\,\big(\psi^\dagger\,\Wg\,\beta\,\ThQ^{K}\psi\big)^{2},
\label{eq:forms}
\end{equation}
to be compared with $I_r^{\mathrm Q}=\sum_{|K|=r}\sgnQ\,(\psi^\dagger\beta\,\ThQ^{K}\psi)^{2}$:
\emph{form C is form Q with the matrix $\Wg$ inserted in every bilinear, and nothing else.} What the
insertion does is seen with the bilinears of the halves (\ref{eq:halves}),
$\mathcal J_K^{++}=\psi_+^\dagger\beta\,\ThQ^{K}\psi_+$,
$\mathcal J_K^{--}=\psi_-^\dagger\beta\,\ThQ^{K}\psi_-$,
$\mathcal J_K^{+-}=\psi_+^\dagger\beta\,\ThQ^{K}\psi_-$ and
$\mathcal J_K^{-+}=(\mathcal J_K^{+-})^\dagger$. If $K$ contains an even number of the indices 5, 6, 7,
the Hermitian matrix $\beta\,\ThQ^{K}$ commutes with $\Wg$ and acts inside each half: form Q contains
$(\mathcal J_K^{++}+\mathcal J_K^{--})^{2}$ and form C contains
$(\mathcal J_K^{++}-\mathcal J_K^{--})^{2}$. The squares, which act between two particles of the same
half, are equal; the product $2\,\mathcal J_K^{++}\mathcal J_K^{--}$, which acts between a particle
of one half and a particle of the other, has opposite signs. If $K$ contains an odd number,
$\beta\,\ThQ^{K}$ anticommutes with $\Wg$ and connects the halves: form Q contains
$(\mathcal J_K^{+-}+\mathcal J_K^{-+})^{2}$ and form C contains
$(\mathcal J_K^{+-}-\mathcal J_K^{-+})^{2}$. The squares $(\mathcal J_K^{+-})^{2}$ and
$(\mathcal J_K^{-+})^{2}$, which move two particles from one half to the other, are equal; the
products of $\mathcal J_K^{+-}$ and $\mathcal J_K^{-+}$, which exchange a particle of one half with a
particle of the other, have opposite signs. \emph{Two particles of the same half interact
identically in the two forms; a particle of the half $\Wg=+1$ and a particle of the half $\Wg=-1$
interact with opposite signs.} Two further consequences are used below. In both forms the
interaction can move two particles together from one half to the other, so that only the total
number of particles is conserved, and the difference of the numbers of the two halves modulo 2. And
both forms are symmetric under the exchange of the two halves, which for fields in ordinary space is
the transformation $\psi\to\gamma^{0}\gamma^{5}\psi$ (a Hermitian matrix with the square 1 that
commutes with $\alpha^{1},\alpha^{2},\alpha^{3}$ and $\beta$ and anticommutes with $\Wg$).

\emph{What the covariant field equation itself says.} If the quartic term of form C is combined with
the covariant kinetic and mass terms (\ref{eq:Lfree}) in place of the positive ones, the field
equation is Eq.~(\ref{eq:dirac}). Restricted to the quantized sector it reads
\begin{equation}
i\partial_4\psi=\hh\,\psi+\Wg\,\frac{\partial\LselfC}{\partial\psi^\dagger},
\label{eq:fieldcov}
\end{equation}
whereas Eq.~(\ref{eq:fieldC}) has $-\partial\LselfC/\partial\psi^\dagger$ in place of
$+\Wg\,\partial\LselfC/\partial\psi^\dagger$. On the half $\Wg=-1$ the two agree. On the half
$\Wg=+1$ the covariant equation has the opposite sign of the whole interaction, the same change of
sign as in Eq.~(\ref{eq:indef}); it conserves $\psi^\dagger\Wg\psi$ rather than $\psi^\dagger\psi$ and
is not the equation of a Hamiltonian with states of positive norm when both halves are occupied. It
is a consistent classical field equation; no quantum theory of particles is built on it here.

Table~\ref{tab:forms} places the three options side by side. The hypotheses fix the nine invariants
as the invariants of the spacetime. On the $7+1$ spacetime this determines the model. On the $4+4$
spacetime it does not, because the adjoint that makes the invariants covariant, $\Cm$, is not the
adjoint that makes the energy positive, $\beta$: the two differ by the matrix $\Wg$. This paper
takes form C as the definition, by the decision of the author of the hypotheses, and reports every
result of form Q beside it.

\begin{table}
\caption{\label{tab:forms}The three ways to complete the self-interaction on the $4+4$ spacetime
in the quantized sector.}
\begin{ruledtabular}
\begin{tabular}{p{0.19\columnwidth}p{0.22\columnwidth}p{0.24\columnwidth}p{0.23\columnwidth}}
option & kinetic and mass term & bilinears of the quartic term & status \\
\hline
form C & positive, Eq.~(\ref{eq:Lpos}) & $\tilde\psi\,\Theta^{K}\psi$, Eq.~(\ref{eq:LselfC}) & the definition \\
form Q & positive, Eq.~(\ref{eq:Lpos}) & $\bar\psi\,\ThQ^{K}\psi$, Eq.~(\ref{eq:LselfQ}) & the alternative, reported beside the definition \\
covariant classical equation & indefinite, Eq.~(\ref{eq:indef}) & as in form C & Eq.~(\ref{eq:fieldcov}); not quantized \\
\end{tabular}
\end{ruledtabular}
\end{table}
%%</1644>

\subsection{The Fierz matrix}

The 256 matrices $\Theta^{K}$ are a basis of all $16\times16$ matrices,
$\mathrm{Tr}(\Theta_K\Theta^{J})=16\,\delta_{KJ}$ for two lists $K$ and $J$, so that
$\sum_K(\Theta_K)_{cd}(\Theta^{K})_{ab}=16\,\delta_{ad}\delta_{bc}$, where the indices
$a,b,c,d=1,\dots,16$ number the components. This completeness relation rearranges a product of two
matrices of the basis, as in the rearrangement of Fierz for four Dirac matrices \cite{Fierz1937}:
\begin{equation}
\sum_{|K|=r}(\Theta_K)_{ab}\,(\Theta^{K})_{cd}=\sum_{r'=0}^{8}\sfF_{rr'}\sum_{|J|=r'}(\Theta_J)_{ad}\,(\Theta^{J})_{cb},
\label{eq:fierz}
\end{equation}
with the \emph{Fierz matrix}
\begin{equation}
\sfF_{rr'}=\frac1{16}\sum_{l}(-1)^{rr'-l}\binom{r'}{l}\binom{8-r'}{r-l} ,
\label{eq:fierzF}
\end{equation}
in which $r$ and $r'$ are ranks and $l$ is a summation index. The derivation takes three lines.
Expanding the left-hand side of Eq.~(\ref{eq:fierz}) in the basis gives the coefficient
$\tfrac1{256}\sum_{|K|=r}\mathrm{Tr}(\Theta_K\Theta_J\Theta^{K}\Theta^{J})$ for each list $J$ of rank
$r'$. Two products of Dirac matrices either commute or anticommute:
$\Theta_K\Theta_J\Theta^{K}=(-1)^{rr'-l}\Theta_J$ if the lists $K$ and $J$ have $l$ indices in common.
And $\binom{r'}{l}\binom{8-r'}{r-l}$ lists $K$ of rank $r$ share exactly $l$ indices with a given $J$
of rank $r'$. Table~\ref{tab:fierz} prints the matrix. It satisfies $\sfF^{2}=1$; its eigenvalues are
$+1$ (five times) and $-1$ (four times).\SIG{ It is the same matrix for both forms (with
$\ThQ^{K}$ for $\Theta^{K}$ in form Q).}{}

\begin{table}
\caption{\label{tab:fierz}The Fierz matrix: the integers $16\,\sfF_{rr'}$ of Eq.~(\ref{eq:fierzF}),
rows $r=0,\dots,8$, columns $r'=0,\dots,8$.}
\begin{ruledtabular}
\begin{tabular}{c|rrrrrrrrr}
$r\backslash r'$ & 0 & 1 & 2 & 3 & 4 & 5 & 6 & 7 & 8 \\
\hline
{{ROWS:tbl_fierz}}
\end{tabular}
\end{ruledtabular}
\end{table}

\subsection{Anticommuting fields: only four invariants are independent}
\label{sec:four}

Contract Eq.~(\ref{eq:fierz}) with two adjoint fields (indices $a$ and $c$) and two fields (indices
$b$ and $d$). The left-hand side is the invariant $I_r$. The right-hand side contains the two fields
$\psi$ in the opposite order, and the components of $\psi$ anticommute, which costs a sign. The nine
invariants therefore obey
\begin{equation}
I_r=-\sum_{r'}\sfF_{rr'}\,I_{r'},\qquad\text{that is}\qquad(1+\sfF)\,I=0 .
\label{eq:theorem}
\end{equation}
Because $\sfF$ has the eigenvalue $+1$ five times, these are five independent linear relations:
\emph{of the nine invariants only four are independent.} Solved for the tensor invariants,
\begin{equation}
\begin{aligned}
I_2&=-3I_0-I_1-I_7-I_8, & I_3&=-4I_0-I_7+4I_8,\\
I_4&=-5I_0-5I_8, & I_5&=-4I_0-I_1+4I_8,\\
I_6&=-I_0+I_1+I_7-3I_8 .
\end{aligned}
\label{eq:reduction}
\end{equation}
Inserted into Eq.~(\ref{eq:Lself}),
\begin{equation}
\Lself=-\tfrac12\big(\geff0\,I_0+\geff1\,I_1+\geff7\,I_7+\geff8\,I_8\big)
\label{eq:Lself4}
\end{equation}
with the four \emph{effective couplings}
\begin{equation}
\begin{aligned}
\geff0&=g_0-3g_2-4g_3-5g_4-4g_5-g_6,\\
\geff1&=g_1-g_2-g_5+g_6,\\
\geff7&=g_7-g_2-g_3+g_6,\\
\geff8&=g_8-g_2+4g_3-5g_4+4g_5-3g_6 .
\end{aligned}
\label{eq:geff}
\end{equation}
A tensor coupling $g_2,\dots,g_6$ is not a new parameter of the theory: it shifts the scalar,
vector, axial-vector and pseudoscalar couplings. Every physical quantity computed below depends on
the nine $g_r$ only through the four combinations (\ref{eq:geff}), and this is verified in each case.
\SIG{The relations (\ref{eq:reduction}) and the combinations (\ref{eq:geff}) are the same in both
forms, with $I_r^{\mathrm C}$ or $I_r^{\mathrm Q}$ for $I_r$: a model is specified, in either form, by
the four numbers $\geff0,\geff1,\geff7,\geff8$, and the choice of the form decides what these four
numbers mean. }{}The theorem (\ref{eq:theorem}) was verified in two independent ways: by exact integer
arithmetic on the $16^{4}$ coefficients of the quartic form, antisymmetrized as anticommuting fields
require, and by evaluating the nine invariants on Grassmann-valued fields with an implementation of
the Grassmann algebra; in both, the rank of the nine invariants is exactly four.

\subsection{Uniform matter: the mean-field couplings}
\label{sec:hf}

%%<1671>
In uniform matter at rest the expectation value of a product of two fields is the density matrix
$\sfn_{ab}=\langle\psi_b^\dagger\psi_a\rangle$. For a filled Fermi sea of particles of mass $\mst$
it is $\sfn=(n+\nS\,\beta)/16$, where $n=\langle\psi^\dagger\psi\rangle$ is the number density and
$\nS=\langle\bar\psi\psi\rangle$ the scalar density, with $\beta=\gamma^{4}$. By Wick's theorem the
expectation value of a normal-ordered product of two bilinears is a direct (Hartree) term minus an
exchange (Fock) term,
$\langle(\bar\psi\Theta^{K}\psi)^{2}\rangle=[\mathrm{Tr}(\beta\Theta^{K}\sfn)]^{2}-\mathrm{Tr}(\beta\Theta^{K}\sfn\,\beta\Theta^{K}\sfn)$.
The energy density of the self-interaction, $-\langle\Lself\rangle$, becomes
\begin{equation}
\varepsilon_{\mathrm{int}}=\tfrac12\,\GV^{\mathrm{HF}}\,n^{2}+\tfrac12\,\GS^{\mathrm{HF}}\,\nS^{2},
\label{eq:EHF}
\end{equation}
\begin{equation}
\begin{aligned}
\GS^{\mathrm{HF}}&=\tfrac1{16}\big(15\,\geff0-8\,\geff1-8\,\geff7-\geff8\big),\\
\GV^{\mathrm{HF}}&=\tfrac1{16}\big(-\geff0+22\,\geff1-6\,\geff7+\geff8\big).
\end{aligned}
\label{eq:GHF}
\end{equation}
These are the \emph{Hartree-Fock couplings}. In terms of the nine couplings, {{GEN:hf_formula_S}}
and {{GEN:hf_formula_V}}. The direct terms alone would give $\GS=g_0$ and $\GV=g_1$; the exchange
terms reduce the scalar coupling by one sixteenth and add to the vector coupling.
Equation~(\ref{eq:EHF}) is the reason why uniform matter is described by two numbers, a scalar
coupling $\GS$ and a vector coupling $\GV$, whatever the nine couplings are: Sec.~\ref{sec:ks}
takes these two as its parameters.
%%</1671>
%%<1644>
In uniform matter at rest the expectation value of a product of two fields is the density matrix
$\sfn_{ab}=\langle\psi_b^\dagger\psi_a\rangle$. When the two halves (\ref{eq:halves}) of the field are
filled with particles up to the Fermi momenta $k_{F+}$ and $k_{F-}$, and no coherence between the
halves is present, it has four terms,
\begin{equation}
\sfn=\tfrac1{16}\big(n+\nS\,\beta+n_W\,\Wg+n_{SW}\,\Wg\beta\big),
\label{eq:nhalves}
\end{equation}
with $n=n_++n_-$, $n_W=n_+-n_-$, $\nS=n_{S+}+n_{S-}$ and $n_{SW}=n_{S+}-n_{S-}$, where
$n_\pm=\langle\psi_\pm^\dagger\psi_\pm\rangle$ are the number densities of the two halves and
$n_{S\pm}=\langle\psi_\pm^\dagger\beta\,\psi_\pm\rangle$ their scalar densities. By Wick's theorem
the expectation value of a normal-ordered product of two bilinears is a direct (Hartree) term minus
an exchange (Fock) term. With it the energy density of the self-interaction is, in both forms, a sum
of four squares,
\begin{equation}
\varepsilon_{\mathrm{int}}=\tfrac12\big(\GVtot\,n^{2}+\GStot\,\nS^{2}+\GVdif\,n_W^{2}+\GSdif\,n_{SW}^{2}\big).
\label{eq:EHF}
\end{equation}
$\GVtot,\GStot$ are the couplings of the squares of the \emph{total} densities, $\GVdif,\GSdif$
those of the squares of the \emph{differences} between the halves. They are four combinations of the
effective couplings,
\begin{equation}
\begin{aligned}
\GS^{\mathrm{HF}}&=\tfrac1{16}\big(15\,\geff0-8\,\geff1-8\,\geff7-\geff8\big),\\
\GV^{\mathrm{HF}}&=\tfrac1{16}\big(-\geff0+22\,\geff1-6\,\geff7+\geff8\big),\\
\GS^{\mathrm c}&=\tfrac1{16}\big(-\geff0-2\,\geff1+2\,\geff7+\geff8\big),\\
\GV^{\mathrm c}&=-\tfrac1{16}\big(\geff0+\geff8\big),
\end{aligned}
\label{eq:GHF}
\end{equation}
\emph{the same four in both forms, with the roles of the totals and of the differences exchanged}:
\begin{equation}
\begin{aligned}
\text{form C:}\quad&\GStot=\GS^{\mathrm c},\quad\GVtot=\GV^{\mathrm c},\\
&\GSdif=\GS^{\mathrm{HF}},\quad\GVdif=\GV^{\mathrm{HF}};\\
\text{form Q:}\quad&\GStot=\GS^{\mathrm{HF}},\quad\GVtot=\GV^{\mathrm{HF}},\\
&\GSdif=\GS^{\mathrm c},\quad\GVdif=\GV^{\mathrm c}.
\end{aligned}
\label{eq:roles}
\end{equation}
In terms of the nine couplings, {{GEN:hf_formula_S}}, {{GEN:hf_formula_V}},
{{GEN:hf_formula_Sc}} and {{GEN:hf_formula_Vc}}. $\GS^{\mathrm{HF}}$ and $\GV^{\mathrm{HF}}$ are the
Hartree-Fock couplings of the $7+1$ model: the direct terms alone would give $g_0$ and $g_1$, and
the exchange terms reduce the scalar coupling by one sixteenth and add to the vector coupling. An
equivalent way to write Eq.~(\ref{eq:EHF}) shows what happens between particles:
\begin{equation}
\begin{aligned}
\varepsilon_{\mathrm{int}}={}&\tfrac12\,\GVsame\,(n_+^{2}+n_-^{2})+\tfrac12\,\GSsame\,(n_{S+}^{2}+n_{S-}^{2})\\
&+\GVopp\,n_+n_-+\GSopp\,n_{S+}n_{S-},
\end{aligned}
\label{eq:Esame}
\end{equation}
with $G^{\mathrm{same}}=G^{\mathrm{tot}}+G^{\mathrm{dif}}$ and
$G^{\mathrm{opp}}=G^{\mathrm{tot}}-G^{\mathrm{dif}}$ for the scalar and for the vector couplings alike:
\begin{equation}
\begin{aligned}
\GSsame&=\tfrac18\big(7\,\geff0-5\,\geff1-3\,\geff7\big),\\
\GVsame&=\tfrac18\big(-\geff0+11\,\geff1-3\,\geff7\big),\\
\GSopp&=\pm\tfrac18\big(8\,\geff0-3\,\geff1-5\,\geff7-\geff8\big),\\
\GVopp&=\pm\tfrac18\big(11\,\geff1-3\,\geff7+\geff8\big),
\end{aligned}
\label{eq:Gsame}
\end{equation}
with the upper sign in form Q and the lower sign in form C. $\GSsame,\GVsame$ are the couplings
between two particles of the same half; they are the same in both forms. $\GSopp,\GVopp$ are the
couplings between a particle of one half and a particle of the other; they have opposite signs in
the two forms. This is the statement of Sec.~\ref{sec:forms} for the mean field.

For an attractive scalar coupling alone ($g_0<0$) the consequence is this. In form Q the attraction
stands in $\GStot=\tfrac{15}{16}g_0$: it lowers the energy of any scalar density, as on $7+1$. In form
C it stands in $\GSdif=\tfrac{15}{16}g_0$: it lowers the energy only if the scalar densities of the two
halves differ, and the total scalar density is weakly repelled ($\GStot=-\tfrac1{16}g_0>0$). Between
two particles of the same half the coupling attracts in both forms ($\GSsame=\tfrac78g_0$); between
particles of opposite halves it attracts in form Q ($\GSopp=g_0$) and repels in form C
($\GSopp=-g_0$). Matter in which only one half is filled is matter of one kind of particle with four
internal states and the couplings $\GSsame$, $\GVsame$; matter of form Q with equal halves is matter
of one kind with eight internal states and the couplings $\GS^{\mathrm{HF}}$, $\GV^{\mathrm{HF}}$.
Section~\ref{sec:ks} sets out the method for one kind of particle with a scalar coupling $\GS$ and a
vector coupling $\GV$, and then treats the two halves together.
%%</1644>

\subsection{Two particles: the pair couplings}
\label{sec:pairs}

Two particles at rest can be in $8\cdot7/2=28$ antisymmetric states, since a particle at rest has
eight internal states (the eigenvalue $+1$ of $\beta$). The self-interaction acts on a pair state
$\Phi$, an antisymmetric $16\times16$ matrix, as
$\Phi\mapsto\sum_rg_r\sum_{|K|=r}\sgnK\,(\SIG{\Cm}{\beta}\Theta^{K})\,\Phi\,(\SIG{\Cm}{\beta}\Theta^{K})^{T}$\SIG{ in form C
(in form Q with $\beta\,\ThQ^{K}$ and $\sgnQ$)}{},
projected on the states of positive energy. Its eigenvalue is the \emph{pair coupling} $\Gpair$,
the interaction energy of the pair times the volume; a negative value means attraction. Diagonalizing
this map (exactly in the program, numerically in the independent verification) gives\SIG{, in form Q
and on the $7+1$ spacetime,}{} two values:
\begin{equation}
\Gpair=\begin{cases}\geff0+\geff1+5\,\geff7&\text{for 7 pairs,}\\ \geff0+\geff1-3\,\geff7&\text{for 21 pairs.}\end{cases}
\label{eq:Gpair}
\end{equation}
The seven pairs are those created by the vector pair operators $\psi^{T}\Cm\gamma^{A}\psi$ along
the seven axes other than the time ($A\neq4$); among them is the pair along the hidden axis $x0$,
the pair of the superfluid of this paper. A scalar attraction ($\geff0<0$) binds all 28 pairs equally; a vector repulsion
($\geff1>0$) opposes all of them equally; the axial coupling distinguishes the seven from the
twenty-one; and the pseudoscalar coupling does not act between particles at rest.
%%<1644>

In form C the result is different. The state that a vector pair operator creates has a definite
behavior under $\Wg$: $\Wg\,\Phi\,\Wg^{T}=-\Phi$ for the pairs along $x0,x1,x2,x3$ and
$\Wg\,\Phi\,\Wg^{T}=+\Phi$ for the pairs along $x5,x6,x7$. \emph{A pair along one of the four
space-like axes consists of one particle of each half; a pair along one of the three additional
time-like axes consists of two particles of the same half.} By Sec.~\ref{sec:forms} the two forms
agree on the pairs of two particles of the same half and have opposite signs on the pairs of one
particle of each half (Table~\ref{tab:pairsrest}). In form C the seven vector pairs split into the
four pairs along the space-like axes, with $\Gpair=-(\geff0+\geff1+5\,\geff7)$, and the three pairs
along the additional axes, with $\Gpair=+(\geff0+\geff1+5\,\geff7)$: \emph{a scalar coupling that
attracts in form Q, and on $7+1$, repels in the hidden-axis pair of form C.}

\begin{table}
\caption{\label{tab:pairsrest}The pair couplings of the 28 antisymmetric states of two particles at
rest on the $4+4$ spacetime, in the two forms of the self-interaction (exact).}
\begin{ruledtabular}
\begin{tabular}{crll}
form & pairs & the two particles are in & $\Gpair$ \\
\hline
{{ROWS:tbl_fm_pairs_rest}}
\end{tabular}
\end{ruledtabular}
\end{table}
%%</1644>

\subsection{Symmetries}
\label{sec:selfsym}

\emph{Chiral symmetry.} The transformation $\psi\to e^{i\epsilon\chi}\psi$ with the chirality matrix
$\chi$ (Hermitian, $\chi^{2}=1$, anticommuting with every $\gamma^{A}$) leaves the self-interaction
invariant if and only if
\begin{equation}
\SIG{\geff1-\geff7=0\ \text{(form C)},\qquad\geff0+\geff8=0\ \text{(form Q)}.}{\geff0+\geff8=0 .}
\label{eq:chiralcond}
\end{equation}
\SIG{In form Q, as on the $7+1$ spacetime, the transformation leaves $I_1$ and $I_7$ unchanged and
rotates $I_0$ and $I_8$ into each other. In form C the roles of the even and of the odd ranks are
exchanged, because the bilinears of the two forms differ by the matrix $\Wg$, which anticommutes with
$\chi$: a scalar coupling alone is chirally invariant in form C and not in form Q, and a vector
coupling alone is chirally invariant in form Q and not in form C. }{The transformation leaves $I_1$
and $I_7$ unchanged and rotates $I_0$ and $I_8$ into each other. }The
mass term breaks this symmetry in any case; it is the symmetry $\UVS$ of the model with the
Vorton-Scalar (Sec.~\ref{sec:clock}), in which the mass arises from the vacuum value of a scalar
field that transforms as well.

\emph{The internal symmetry of fields in ordinary space.} A field that does not depend on the four
extra coordinates $x0,x5,x6,x7$ has an internal symmetry $[SU(2)\times SU(2)]\times U(1)_E$ in
addition to the phase symmetry $U(1)_F$ of the fermion number (Sec.~\ref{sec:higgs}): six generators
proportional to $\gamma^{a}\gamma^{b}$ ($a<b$ out of $0,5,6,7$), which generate the rotations $SO(4)$
of the four extra axes, and a generator $P$ proportional to $\gamma^{0}\gamma^{5}\gamma^{6}\gamma^{7}$.
%%<1671>
The self-interaction is invariant under the six rotations of the extra axes and under the phase for
all couplings, but under $U(1)_E$ only if $\geff1=\geff7=0$: \emph{a vector or an axial-vector contact
coupling breaks $U(1)_E$.} A scalar and a pseudoscalar coupling keep it.
%%</1671>
%%<1644>
The three rotations of $x5,x6,x7$ among themselves commute with $\Wg$; the three generators
$\gamma^{0}\gamma^{t}$ ($t=5,6,7$) and $P$ anticommute with $\Wg$ and connect the two halves.
Table~\ref{tab:intsym} gives the couplings for which the self-interaction is invariant under the
unitary transformations that each generator generates. \emph{In form C every self-interaction that
does not vanish breaks the three transformations that mix the hidden axis with $x5,x6,x7$, and
breaks $U(1)_E$}: the internal symmetry is reduced from $SO(4)\times U(1)_F\times U(1)_E$ to
$SO(3)_T\times U(1)_F$, the rotations of the three additional times among themselves and the fermion
number. In form Q, as on $7+1$, the self-interaction keeps the six rotations and the phase for all
couplings and keeps $U(1)_E$ only if $\geff1=\geff7=0$. The reason for the difference is this. In the
frame group of the $4+4$ spacetime $\gamma^{0}\gamma^{t}$ generates a boost, which mixes a space-like
with a time-like coordinate and is not a unitary transformation of the field. The unitary
transformation generated by the same matrix is a different transformation: it is a symmetry of the
kinetic term of fields in ordinary space by accident, and it is not an element of the frame group.
The quartic term of form C is invariant under the boost and has no reason to be invariant under the
unitary transformation. In form Q the same matrices generate rotations, as on the $7+1$ spacetime,
where all four extra axes are space-like.

\begin{table*}
\caption{\label{tab:intsym}The couplings for which the self-interaction is invariant under the
internal symmetry of fields in ordinary space, on the $7+1$ spacetime and in the two forms of the
$4+4$ spacetime (exact).}
\begin{ruledtabular}
\begin{tabular}{lcccc}
generators & number & $7+1$ & $4+4$, form Q & $4+4$, form C \\
\hline
{{ROWS:tbl_fm_symmetry}}
\end{tabular}
\end{ruledtabular}
\end{table*}
%%</1644>
